\chapter{KẾT LUẬN VÀ HƯỚNG PHÁT TRIỂN}

\section{Kết luận}

\subsection{Kết quả đạt được}

Đồ án đã hoàn thành việc nghiên cứu, thiết kế và triển khai thực tế hệ thống dịch tài liệu đa
ngôn ngữ cho môi trường doanh nghiệp. Các kết quả cụ thể đạt được bao gồm:

\begin{itemize}
  \item \textbf{Hệ thống dịch thuật đa định dạng:} Hỗ trợ dịch sáu định dạng tài liệu
    (DOCX, XLSX, PPTX, PDF native-text, PDF scan, CSV/XLS) sang tối đa 10 ngôn ngữ
    đồng thời trong một yêu cầu. Định dạng và bố cục tài liệu được bảo toàn sau dịch.

  \item \textbf{Pipeline PDF thông minh:} Thuật toán phân loại PDF ba tiêu chí
    (\texttt{is\_pdf\_truly\_editable}) hoạt động chính xác, tự động định tuyến tài liệu đến
    pipeline ConvertAPI→DOCX hoặc Azure OCR phù hợp mà không cần can thiệp của người dùng.

  \item \textbf{Hệ thống quản lý thuật ngữ:} Ba cấp (Private Library, Suggestion Queue,
    Common Library) với 44 cặp glossary Google Translation API hoạt động tự động, cho phép
    doanh nghiệp duy trì bộ từ điển kỹ thuật nhất quán qua cơ chế đề xuất cộng đồng.

  \item \textbf{Kiến trúc sạch:} Cơ chế thread-local context injection cho phép tích hợp chế
    độ thư viện vào tất cả module dịch mà không thay đổi chữ ký hàm hiện có, đảm bảo tính
    module hóa của codebase.

  \item \textbf{Triển khai production:} Hệ thống đã được triển khai tại
    AWS với Cognito SSO, AWS~S3, Application Load Balancer và container Docker, phục vụ
    nhu cầu dịch thuật tài liệu kỹ thuật đa ngôn ngữ trong môi trường sản xuất thực tiễn.
\end{itemize}

\subsection{Hạn chế}

\begin{itemize}
  \item \textbf{Thời gian xử lý tài liệu lớn:} Tài liệu nhiều trang (>~50 trang) có thể mất
    vài phút do tính chất tuần tự của các lần gọi Google Translation~API và pipeline OCR.
    Hiện chưa có cơ chế ước lượng tiến độ (progress bar) trên giao diện.

  \item \textbf{Bảo toàn định dạng PPTX phức tạp:} Một số slide PPTX có bố cục phức tạp
    (SmartArt, grouping) có thể bị mất định dạng sau dịch do giới hạn của python-pptx khi
    xử lý các đối tượng đặc biệt.

  \item \textbf{Cập nhật glossary bất đồng bộ:} Sau khi duyệt thuật ngữ, glossary Google
    Translation được cập nhật trong background thread (30--180~giây/pair × 44 pairs). Trong
    khoảng thời gian này, các yêu cầu dịch vẫn sử dụng glossary cũ.

  \item \textbf{Phát hiện ngôn ngữ nguồn:} Một số tài liệu hỗn hợp nhiều ngôn ngữ trong cùng
    một tệp có thể bị phát hiện sai ngôn ngữ nguồn, ảnh hưởng đến việc chọn glossary phù hợp.
\end{itemize}

\section{Hướng phát triển}

Dựa trên kết quả đạt được và các hạn chế còn tồn tại, đồ án đề xuất các hướng phát triển trong
tương lai:

\begin{enumerate}
  \item \textbf{Xử lý bất đồng bộ với Celery/Redis:} Chuyển pipeline dịch tài liệu sang hệ
    thống hàng chờ tác vụ (Celery + Redis/RabbitMQ) để hỗ trợ giao diện theo dõi tiến độ
    thời gian thực (WebSocket) và cải thiện khả năng mở rộng.

  \item \textbf{Gợi ý thuật ngữ bằng ML:} Áp dụng mô hình NLP để tự động gợi ý các cặp
    thuật ngữ chưa có trong thư viện dựa trên nội dung tài liệu đã dịch, hỗ trợ người dùng
    xây dựng thư viện cá nhân nhanh hơn.

  \item \textbf{Ứng dụng di động:} Phát triển ứng dụng mobile (React Native hoặc Flutter) cho
    phép nhân viên nhà máy chụp ảnh tài liệu và dịch ngay trên thiết bị di động, tích hợp
    với OCR pipeline hiện có.

  \item \textbf{Dashboard phân tích nâng cao:} Mở rộng trang thống kê Admin với biểu đồ xu
    hướng sử dụng theo thời gian, top ngôn ngữ dịch, tài liệu dịch nhiều nhất, giúp quản lý
    đưa ra quyết định về ưu tiên mở rộng ngôn ngữ.

  \item \textbf{Cache kết quả dịch:} Triển khai Redis cache cho các đoạn văn bản lặp lại
    (đặc biệt trong tài liệu kỹ thuật có nhiều đầu mục cố định), giảm số lần gọi API và
    chi phí vận hành.

  \item \textbf{Hỗ trợ dịch ngoại tuyến:} Tích hợp mô hình dịch máy cục bộ (ví dụ:
    Helsinki-NLP/Opus-MT qua HuggingFace) cho các tài liệu nội bộ nhạy cảm không muốn
    gửi ra dịch vụ đám mây bên ngoài.
\end{enumerate}
